% =====================================================================
\section{Mathematical results}\label{sec:math}
% =====================================================================

Throughout this section we work under the following standing
hypotheses.

\begin{assumption}[Standing hypotheses]\label{ass:standing}
The triple $(\Q,E,\Psi)$ with $E(\bs q,S)=W(\bs q)-\bs q\cdot\bs\ell(S)$
satisfies:
\begin{enumerate}[label=\textnormal{(A\arabic*)},leftmargin=2.6em,itemsep=2pt]
  \item\label{A:W} \emph{Stored energy.} $\Q\subseteq\R^n$ is closed and
  $W\colon\Q\to\R$ is continuous and coercive of order $p>1$: there
  exist $c_1>0$, $c_2\ge0$ with $W(\bs q)\ge c_1\|\bs q\|^p-c_2$.
  \item\label{A:ell} \emph{Interaction potential.}
  $\bs\ell\colon[0,\infty)\to\R^n$ is Lipschitz and bounded, with
  $C_\ell:=\sup_{S\ge0}\|\bs\ell(S)\|<\infty$ and
  $C_{\ell'}:=\sup_{S\ge0}\|\bs\ell'(S)\|<\infty$.
  \item\label{A:S} \emph{Loading.} $S\in W^{1,1}(0,T;\R^+)$.
  \item\label{A:Psi} \emph{Dissipation potential.}
  $\Psi\colon\R^n\to[0,\infty)$ is continuous, convex, normalised
  $(\Psi(\bs 0)=0)$, positively $1$-homogeneous and symmetric
  $(\Psi(-\bs v)=\Psi(\bs v))$.
  \item\label{A:coerc} \emph{Definiteness of dissipation.} There
  exists $c_\Psi>0$ with $\Psi(\bs v)\ge c_\Psi\|\bs v\|$ for all
  $\bs v$; equivalently, $\bs 0\in\operatorname{int}\E$.
\end{enumerate}
\end{assumption}

Hypothesis~\ref{A:coerc} is used only for the total-variation bound in
the balanced-viscosity theory (\cref{sec:bv}); it holds with
$c_\Psi=\rho$ for $\Psi(\bs v)=\rho\|\bs v\|$ and, more generally,
whenever the elastic domain has the origin in its interior.

\subsection{Properties of the dissipation distance}\label{sec:diss_props}

\begin{proposition}[Quasi-pseudodistance]\label{prop:quasi_distance}
Under~\ref{A:Psi}, the dissipation distance~\eqref{eq:diss_dist}
satisfies, for all $\bs q_1,\bs q_2,\bs q_3\in\Q$,
\[
  \Dr(\bs q_1,\bs q_2)\ge0,\quad \Dr(\bs q,\bs q)=0,\qquad
  \Dr(\bs q_1,\bs q_3)\le\Dr(\bs q_1,\bs q_2)+\Dr(\bs q_2,\bs q_3).
\]
If $\Psi$ is symmetric then $\Dr$ is a pseudodistance; if in addition
$\Psi(\bs v)=0\Rightarrow\bs v=\bs 0$ (e.g.\ under~\ref{A:coerc}), then
$\Dr$ is a distance.
\end{proposition}

\begin{proposition}[Linearity and continuity]\label{prop:Dr_linear}
If $\Q$ is convex and $\Psi$ is continuous, convex and positively
$1$-homogeneous, then
\begin{equation}\label{eq:Dr_linear}
  \Dr(\bs q_1,\bs q_2)=\Psi(\bs q_2-\bs q_1)\qquad\forall\,\bs q_1,\bs q_2\in\Q,
\end{equation}
and $\Dr$ is continuous on $\Q\times\Q$. In particular, for
$\Psi(\dot q)=\rho|\dot q|$ one has $\Dr(q_1,q_2)=\rho|q_2-q_1|$.
\end{proposition}

The identity~\eqref{eq:Dr_linear} is central to the numerical method:
it makes the incremental dissipation independent of the step size.

\subsection{Energy estimates}\label{sec:energy_estimates}

\begin{proposition}[Coercivity, absolute continuity and power
control]\label{prop:power}
Under~\ref{A:W}--\ref{A:S}, the following hold.
\begin{enumerate}[label=\textnormal{(E\arabic*)},leftmargin=2.6em,itemsep=2pt]
  \item\label{E:compact} For every $c\in\R$ and $S\ge0$ the sublevel
  set $\{\bs q\in\Q:E(\bs q,S)\le c\}$ is compact.
  \item\label{E:AC} For each $\bs q\in\Q$, $t\mapsto E(\bs q,S(t))$ is
  absolutely continuous with
  $\frac{\dd}{\dd t}E(\bs q,S(t))=-\bs q\cdot\bs\ell'(S(t))\dot S(t)$
  a.e.
  \item\label{E:gronwall} There exist $c_3,c_4>0$, depending only on
  the data, with $|\partial_S E(\bs q,S)|\le c_3\bigl(E(\bs q,S)+c_4\bigr)$
  for all $\bs q\in\Q$, $S\ge0$.
\end{enumerate}
\end{proposition}

\subsection{Existence of energetic solutions}\label{sec:existence}

We first recall the set of states that pass the stability test.

\begin{definition}[Stability set]\label{def:stability_set}
For $t\in[0,T]$, the \emph{stability set} is
$\mathcal S(t):=\{\bs q\in\Q: E(\bs q,S(t))\le E(\tilde{\bs q},S(t))+\Dr(\bs q,\tilde{\bs q})\ \ \forall\,\tilde{\bs q}\in\Q\}$.
\end{definition}

The passage to the limit in the existence proof rests on the following
compatibility properties, which we state explicitly (the corresponding
statement is only sketched in the general theory).

\begin{lemma}[Compatibility]\label{lem:compatibility}
Under~\ref{A:W}--\ref{A:Psi}, let $(t_m,\bs q_m)$ be a stable sequence,
i.e.\ $\bs q_m\in\mathcal S(t_m)$ with
$\sup_m E(\bs q_m,S(t_m))<\infty$, and suppose $t_m\to t$ and
$\bs q_m\to\bs q$ in $\Q$. Then:
\textnormal{(C1)} the power converges, $\partial_S E(\bs q_m,S(t_m))\to
\partial_S E(\bs q,S(t))$ for a.e.\ $t$; and \textnormal{(C2)} the
stability set is closed, $\bs q\in\mathcal S(t)$.
\end{lemma}

\begin{theorem}[Existence of energetic solutions]\label{thm:existence}
Under~\ref{A:W}--\ref{A:Psi}, for every initial datum
$\bs q_0\in\mathcal S(0)$ there exists at least one energetic solution
$\bs q\colon[0,T]\to\Q$ of bounded variation, in the sense of
\cref{def:energetic}.
\end{theorem}

The proof verifies the abstract hypotheses of
\cite[Thm.~2.1.6]{MielkeRoubicek} using
\cref{prop:quasi_distance,prop:power,lem:compatibility}; see
\cref{app:math}. For the benchmark of \cref{sec:experiments}, $\Q=\R$
is not compact but $W$ is a coercive quartic, so~\ref{A:W} holds and
\cref{thm:existence} applies.

\subsection{Uniqueness under uniform convexity}\label{sec:uniqueness}

Energetic solutions need not be unique: at a jump, several post-jump
states may satisfy~(S)--(E). Uniqueness is recovered under convexity.

\begin{theorem}[Uniqueness]\label{thm:uniqueness}
Assume, in addition to~\ref{A:W}--\ref{A:Psi}, that
$W\in C^3(\Q)$ is uniformly convex ($D^2W\succeq\alpha\Id$, $\alpha>0$)
and that $\bs\ell\in C^3$, $S\in C^3([0,T])$. Then for every
$\bs q_0\in\mathcal S(0)$ the energetic solution is unique.
\end{theorem}

Uniform convexity forces a single well and thus excludes the
multi-well energies of interest; the scalar theory of \cref{sec:bv}
provides an alternative route.

\subsection{Balanced-viscosity solutions}\label{sec:bv}

\paragraph{Motivation.}
At a jump point $t^\ast$, conditions~(S)--(E) require the released
energy to cover the dissipation cost
$\Dr(\bs q^-(t^\ast),\bs q^+(t^\ast))$ but do not select the post-jump
state. Balanced-viscosity (BV) solutions~\cite{mielke2012bv} resolve
this by selecting the jump path as the vanishing-viscosity limit
($\varepsilon\to0^+$) of the regularised inclusion
\begin{equation}\label{eq:regularised}
  \bs 0\in\partial\Psi(\dot{\bs q}_\varepsilon)
  +\varepsilon\,\dot{\bs q}_\varepsilon
  +\frac{\partial E}{\partial\bs q}(\bs q_\varepsilon,S),
\end{equation}
whose memory is retained in a correction concentrated at jumps.

\begin{definition}[Balanced-viscosity solution]\label{def:BV}
A function $\bs q\colon[0,T]\to\Q$ of bounded variation is a
\emph{BV solution} of $(\Q,E,\Psi)$ if:

\smallskip
\noindent\textbf{(BV-S) Local stability.} For a.e.\ $t$, there is
$\delta>0$ such that $E(\bs q(t),S(t))\le E(\bs r,S(t))+\Dr(\bs q(t),\bs r)$
for all $\bs r\in B_\delta(\bs q(t))\cap\Q$.

\smallskip
\noindent\textbf{(BV-E) Energy balance with viscous correction.}
For every $t$,
\begin{equation}\label{eq:BV_energy}
  E(\bs q(t),S(t))+\Diss_\Psi(\bs q;[0,t])+\mathcal V(\bs q;[0,t])
  = E(\bs q_0,S(0))+\int_0^t\partial_S E(\bs q,S)\dot S\dd s,
\end{equation}
where the \emph{viscous correction} is
\begin{equation}\label{eq:viscous_correction}
  \mathcal V(\bs q;[0,t]):=\sum_{s\le t}
  \bigl[E(\bs q^-(s),S(s))-E(\bs q^+(s),S(s))-\Dr(\bs q^-(s),\bs q^+(s))\bigr]^+,
\end{equation}
the sum ranging over the (at most countably many) jump points.
\end{definition}

\begin{remark}[Jump condition]\label{rem:jump}
$\mathcal V$ is fully determined by the jump structure: at each jump
$t^\ast$ the summand encodes the \emph{jump condition}
$E(\bs q^-(t^\ast),S(t^\ast))-E(\bs q^+(t^\ast),S(t^\ast))
=\Dr(\bs q^-(t^\ast),\bs q^+(t^\ast))+\text{(viscous excess)}$,
selecting the path of steepest descent through the energy barrier.
\end{remark}

\begin{theorem}[Existence of BV solutions]\label{thm:BV_existence}
Under~\ref{A:W}--\ref{A:coerc}, for every $\bs q_0\in\mathcal S(0)$
there exists at least one BV solution in the sense of \cref{def:BV}.
\end{theorem}

\begin{theorem}[Scalar uniqueness]\label{thm:BV_uniqueness}
Let $n=1$, $\Q\subseteq\R$, and assume, besides
\ref{A:W}--\ref{A:coerc}, that $q\mapsto E(q,S)$ has at most finitely
many local minima for each $S\ge0$. Then the BV solution is unique.
\end{theorem}

\cref{thm:BV_uniqueness} covers the scalar double-well energy of the
benchmark and of the biological application, for which
\cref{thm:uniqueness} fails. Since every energetic solution is a BV
solution, uniqueness of BV solutions implies that the energetic
solution, when it is BV, is the unique one. The proofs of the results
of this section are collected in \cref{app:math}.
